% Modernised reading — fonds Alexandre Grothendieck, shelfmark 29, batch 1
% (pages 1-20).
%
% This is an INTERPRETATION, not a transcription. It re-reads the pages in
% today's notation and vocabulary, states the hypotheses the manuscript leaves
% implicit, and drops the critical apparatus so that the mathematics reads
% straight through. Everything the brackets and underlines carried — a doubtful
% reading, a notation he used differently, a missing passage — moves into a
% footnote.
%
% It is held to being mathematically correct as it stands, not merely faithful
% to the page. Where the manuscript is loose or mistaken, the true statement is
% what appears here, and the footnote says what the page has. In any dispute of
% substance the transcription governs: whoever cites Grothendieck must cite
% batch-01.fr.tex.
%
% Scope: the folder was read WHOLE before any of this was written — all eleven
% batches, pages 1-216, in one pass — and one file was then written per batch,
% which is the layout the site reads. What this file knew when it was written
% is the whole shelfmark, and it uses that: the divisibility correction below
% is settled by pages 2 and 37, and the notation is the folder's, not this
% batch's.
%
% Notation fixed here and held for the eleven files: Et(S) for the étale site,
% Fl_S for the fibred category of relative schemes, calligraphic R = (C, s) for
% the 1967 *donnée de ramification*, Rev^R(S) always with the exponent (the
% typescript writes both Rev^R and Rev_R), r for the geometric realisation.
% The roman (R, r) of pages 190-216 is a DIFFERENT object — the champ itself,
% not an auxiliary datum — and batch 10's reading says so at the point where
% the two meet.
%
% Three departures from the page, each footnoted where it occurs:
%   - page 20's « n' >= n » is the divisibility order, not the product order:
%     Murre queries it on page 2 and Grothendieck confirms it on page 37;
%   - page 13's full faithfulness of the transition functors is a standing
%     HYPOTHESIS on the examples, not a proposition;
%   - page 18 writes f : X -> X' where its own objects are Z, Z'.
%
% Pass: Opus 5 (claude-opus-5), 2026-08-29 — first pass, unchecked against the
% pages by a human.
\documentclass[11pt,a4paper]{article}
% Le préambule commun à toutes les transcriptions du fonds Grothendieck.
%
% Il tient en une idée : **l'incertitude se compose, elle ne se résout pas.**
% Une transcription qui lisse un mot douteux a détruit la seule chose pour
% laquelle on la faisait — un lecteur doit pouvoir savoir, à chaque endroit,
% ce qui était sur le feuillet et ce qui a été ajouté. Les six macros qui
% suivent sont donc l'appareil critique, et non de la décoration.
%
% Les mêmes commandes sont reconnues par `scripts/render.mjs`, qui produit la
% vue de lecture du site. Une macro ajoutée ici doit l'être aussi là-bas,
% faute de quoi le rendu échouera bruyamment — ce qui est voulu : un rendu
% silencieusement faux est pire qu'un rendu absent.

\ProvidesPackage{grothendieck}[2026/08/08 Transcriptions du fonds Grothendieck]

\usepackage{amsmath,amssymb,amsthm}
\usepackage{mathrsfs}
\usepackage{stmaryrd}
% Les diagrammes commutatifs sont partout dans ce fonds : c'est la langue
% même de la géométrie algébrique de Grothendieck, et une transcription qui
% les rendrait en prose ne transcrirait rien.
\usepackage{tikz-cd}
% Français par défaut — c'est la langue des feuillets — mais l'anglais est
% chargé aussi : la traduction et le résumé sont des documents anglais, et la
% typographie française (espaces avant les deux-points) y serait une faute.
% Ces documents basculent par \selectlanguage{english} en tête de corps.
\usepackage[english,french]{babel}
\usepackage{fontspec}
\usepackage{geometry}
\geometry{a4paper,margin=25mm,marginparwidth=28mm}
\usepackage{marginnote}
\usepackage{xcolor}
% ulem, not soul. `soul`'s \st and \ul are built on a character-by-character
% scan that XeLaTeX's font machinery breaks: the document fails with an
% undefined control sequence rather than degrading. ulem does the same two jobs
% and is Unicode-engine safe. `normalem` stops it hijacking \emph, which the
% transcriptions use for Grothendieck's own emphasis.
\usepackage[normalem]{ulem}
\usepackage{hyperref}
\hypersetup{colorlinks=true,linkcolor=black,urlcolor=black,pdfborder={0 0 0}}

% --- Métadonnées du lot -----------------------------------------------------
% Reprises telles quelles par le script de rendu, qui les affiche en tête de
% la vue de lecture. Elles fixent ce que le fichier prétend couvrir : sans
% elles, une transcription flotte sans point d'attache dans un fonds de seize
% mille feuillets.

\newcommand{\folder}[1]{\gdef\grothFolder{#1}}
\newcommand{\batch}[1]{\gdef\grothBatch{#1}}
\newcommand{\leaves}[2]{\gdef\grothFirst{#1}\gdef\grothLast{#2}}
\newcommand{\foldertitle}[1]{\gdef\grothFoldertitle{#1}}
\gdef\grothFolder{?}\gdef\grothBatch{?}\gdef\grothFirst{?}\gdef\grothLast{?}\gdef\grothFoldertitle{}

% --- Le résumé --------------------------------------------------------------
%
% Il ouvre la lecture modernisée, sous le titre « Résumé », et il est composé
% en petit. Ce n'est pas un ornement : c'est le seul passage du document qui
% ne soit pas des mathématiques, et un lecteur doit pouvoir le reconnaître
% avant de l'avoir lu — puis le sauter s'il connaît déjà le sujet, ou s'y
% arrêter s'il ne le connaît pas. Le filet qui le suit marque la reprise du
% corps.

\newenvironment{resume}
  {\small\setlength{\parskip}{0.45em}}
  {\par\medskip\hrule\medskip}

% --- Mots-clés ---------------------------------------------------------------
%
% En anglais, en clôture du résumé de la lecture modernisée : le vocabulaire
% moderne sous lequel chercher ce que les feuillets construisent. Le manifeste
% du site (`npm run manifest`) les extrait pour étiqueter la cote entière —
% cette ligne est la source unique des tags, il n'y a pas de fichier à côté.
\newcommand{\keywords}[1]{\par\smallskip\noindent
  {\footnotesize\textsc{Keywords} --- \textit{#1}}\par}

% --- L'appareil critique ----------------------------------------------------

% Le numéro de feuillet, en marge. C'est l'ancre qui permet de régler un
% désaccord sur une formule en regardant *un* feuillet plutôt qu'une chemise
% de six cents. Le site s'en sert aussi pour tourner les pages du fac-similé
% quand on fait défiler la transcription.
\newcommand{\leaf}[1]{%
  \marginnote{\footnotesize\color{gray!70}\textsf{#1}}%
  \label{leaf:#1}\ignorespaces}

% Illisible : on ne devine pas. Un mot inventé qui se lit comme les autres est
% le pire résultat possible.
\newcommand{\ill}{\textcolor{red!60!black}{[\,\ldots\,]}}

% Lecture douteuse : le mot est donné, et signalé comme incertain.
\newcommand{\uncertain}[1]{\uline{#1}}

% Ajout de l'éditeur — une lettre manquante, un mot sous-entendu.
\newcommand{\add}[1]{\textcolor{blue!50!black}{[#1]}}

% Biffé par Grothendieck lui-même. Ce qu'il a rayé fait partie du document :
% une rature dit souvent où la pensée a changé de direction.
\newcommand{\struck}[1]{\sout{#1}}

% Note du transcripteur, sur l'état matériel du feuillet.
\newcommand{\note}[1]{\par\smallskip{\small\color{gray!60!black}\textit{[#1]}}\par\smallskip}

% Note marginale de Grothendieck, distincte du corps du texte.
\newcommand{\marginal}[1]{\par\smallskip\noindent\hspace{1em}%
  {\small\color{gray!60!black}#1}\par\smallskip}

% --- Titre ------------------------------------------------------------------

\AtBeginDocument{%
  \begin{center}
    {\large\bfseries Fonds Alexandre Grothendieck --- cote n\textsuperscript{o}~\grothFolder}\\[2pt]
    {\itshape\grothFoldertitle}\\[4pt]
    {\small feuillets \grothFirst--\grothLast\ (lot \grothBatch)}\\[2pt]
    {\footnotesize\color{gray!60!black}Transcription établie d'après le fac-similé numérisé
     par l'Université de Montpellier.\\
     Les lectures incertaines sont soulignées, les illisibles notées [\,\ldots\,],
     les ajouts entre crochets.}
  \end{center}
  \vspace{1em}}

\endinput


\folder{29}
\batch{1}
\pages{1}{20}
\dating{1959-1969}
\watermark{Édition de démonstration\\interprétation personnelle de l'œuvre}
\foldertitle{Groupe fondamental [Autour de SGA 4 et SGA 7] — lecture modernisée}

\begin{document}

\section*{Résumé}

\begin{resume}
Ces pages ouvrent un dossier de deux cent seize feuillets qui, de 1959 à 1969,
tourne autour d'une seule question : \emph{qu'est-ce qu'un revêtement, en
géométrie algébrique, quand on autorise la ramification} ?

Le mot revêtement vient de la topologie. Au-dessus d'un espace, on en dispose
un autre, qui le recouvre en plusieurs exemplaires, sans jamais se déchirer ni
se coller : un escalier en colimaçon au-dessus d'une cour circulaire, qui monte
d'un étage à chaque tour. Toute l'information est dans la façon dont les
étages se permutent quand on fait le tour ; c'est ce qu'on appelle la
monodromie, et le groupe fondamental n'est rien d'autre que le registre de
toutes les monodromies possibles.

La difficulté commence quand on bouche le trou. Rebouchez le puits central de
la cour : l'escalier, prolongé, retombe sur un simple plancher, et le nombre
de tours n'est plus lisible sur le résultat. C'est exactement la situation d'un
revêtement ramifié : deux données très différentes ont la même
\emph{réalisation géométrique}, et cette réalisation ne les détermine pas. Un
revêtement modérément ramifié n'est donc pas un revêtement qui satisfait une
\emph{propriété} : c'est un revêtement muni d'une \emph{structure}
supplémentaire. Les onze premières pages du dossier sont le procès-verbal de
cette découverte, faite à deux : Jean-Pierre Murre, qui rédigeait pour SGA 1
l'exposé sur le groupe fondamental modéré, écrit qu'il ne voit pas d'où la
structure pourrait venir ; Grothendieck lui répond que ses propres notes
étaient fautives, énumère cinq énoncés faux, et propose de tout reprendre.

Ce que la reprise apporte est le geste central du dossier, et il vaut d'être
énoncé sans jargon : \emph{cesser d'étudier l'objet, étudier la catégorie des
objets}. Un revêtement isolé ne dit pas de quel type de ramification il
provient ; mais la collection de tous les revêtements d'un type donné, avec ses
sommes, ses produits et ses quotients, en dit assez pour reconstituer un
groupe. C'est le principe de la théorie de Galois — un corps ne se connaît pas
tout seul, il se connaît par ses extensions — transposé aux schémas, puis
poussé jusqu'au bout : la « donnée de ramification » que ces pages définissent
n'est qu'un échafaudage, et le dossier passera dix ans à l'ôter.

Il l'ôtera pour de bon : ses dernières pages remplacent l'échafaudage par un
\emph{champ de catégories multigaloisiennes} au-dessus du site étale, ce que
Grothendieck appelle un domaine de ramification, et dont les exemples
kummériens sont ce qu'on nomme aujourd'hui des champs de racines. Entre les
deux, le dossier fait autre chose encore : une lettre à Serre du 18 octobre
1959, au milieu du carton, définit le groupe fondamental des revêtements
\emph{infinitésimaux} — un groupe qui n'a qu'un point et qui n'est pas trivial
pour autant.

L'ordre du carton n'est pas celui de la pensée : 1959 est à la page 124, 1967
au début, et la synthèse à la fin sans date. Les noms modernes à chercher
ensuite sont \emph{groupe fondamental modéré}, \emph{champ}, \emph{champ des
racines}, et \emph{géométrie logarithmique}.

\keywords{tame fundamental group, ramification datum, fibred category, Kummer covering, cartesian section}
\end{resume}

\pagerange{1}{11}
\section*{Ce que Murre n'arrivait pas à voir (pages 1 à 11)}

La chemise porte, de la main de Grothendieck, deux mots : « Tame
ramification ». Dessous, quatre lettres de 1967, que le carton range dans le
désordre. Remises dans l'ordre où elles ont été écrites, elles racontent une
histoire complète.

\subsection*{29 mars 1967 : Murre envoie le manuscrit}

Murre rédige, pour le séminaire de 1960--61, l'exposé sur le groupe fondamental
modéré. Il envoie son manuscrit et sa propre liste de doutes : il a supposé
plus tôt que nécessaire l'existence des quotients ; il hésite sur le mot
\emph{covering}, qui traduit \emph{revêtement} mais entre en collision avec le
« recouvrement » d'une topologie ; il n'a pas trouvé de démonstration plus
simple pour deux de ses lemmes ; et il juge insatisfaisant, sans savoir
l'amender, l'énoncé de son théorème 9.4.

\subsection*{La réponse : cinq énoncés faux, et un diagnostic}

Grothendieck répond que les erreurs sont les siennes --- ses notes étaient trop
sommaires --- et il en donne la liste. La première est un contre-exemple, et
c'est le seul endroit du lot où il calcule.

Soit $S$ un schéma, $b$ une section de $\mathcal{O}_{S}$, $n \geqslant 1$
inversible sur $S$, et posons $a = b^{n}$. Soit
\[
  Y' \;=\; \operatorname{Spec} \mathcal{O}_{S}[T]/(T^{n} - a) ,
\]
revêtement kummérien élémentaire, principal sous $H' = \mu_{n}$. Prenons pour
$H$ le groupe unité et $Y = S$.\footnote{Le tapuscrit écrit $Y\,T/(T^{n}-a)$ :
la machine ne portait pas de crochets, et $Y[T]/(T^{n}-a)$ est ce qu'il faut
lire. La transcription le signale ; c'est le premier des blancs mécaniques de
ce carbone, dont les flèches et les lettres grecques sont dans le même cas.}
Un morphisme de $Y$ dans $Y'$ compatible avec l'inclusion $H \to H'$ est
simplement une section de $Y'$ au-dessus de $S$ : elle existe, $T \mapsto b$.
Or $H \to H'$ n'est pas surjectif dès que $n > 1$. L'énoncé 1.16 des notes de
Murre --- qu'un tel morphisme force $H \to H'$ à être surjectif --- est donc
faux, et avec lui 3.6, puis les démonstrations de 3.7 et 3.8 qui s'en
servaient.

Le théorème 6.4 (le théorème d'Abhyankar) n'est correct, sous la forme donnée,
que si les diviseurs $D_{i}$ sont réguliers ; Grothendieck avertit au passage
que \emph{les groupes d'inertie ne sont déterminés qu'à automorphisme intérieur
près}, ce qui est la raison de la panne, et propose de remplacer la famille
$(D_{i})$ par un unique diviseur $D$ à croisements normaux, en demandant que
les revêtements soient modérément ramifiés localement pour la topologie étale
le long des composantes irréductibles locales de $D$. C'est la notion qu'il
appellera page 21 « modérément ramifié \emph{au sens large} », et c'est elle,
dit-il, qui convient au paragraphe 9.

La démonstration de 7.1 achoppe, elle, sur un point de principe : même pour
$D = 0$ elle aurait à faire intervenir des revêtements étales connexes
\emph{non galoisiens}. D'où le diagnostic, qui est l'origine de tout le reste :
\emph{il faut une notion de ramification modérée pour les revêtements non
galoisiens}, et les hypothèses de normalité et de réduction dont les notes se
servent sont artificielles.

\subsection*{29 avril 1967 : le plan en six paragraphes}

La lettre suivante accompagne l'esquisse que reproduisent les pages 12 à 20, et
propose de refondre les cinq premiers paragraphes de l'exposé en six :
revêtements kummériens relatifs à une famille de sections ; données de
ramification ; propriétés de la catégorie des revêtements de type
$\mathcal{R}$ et le groupe $\pi_{1}^{\mathcal{R}}(S,\;)$ ; fonctorialités ;
revêtements galoisiens de type $\mathcal{R}$ ; théorème d'Abhyankar pour les
diviseurs à croisements normaux.

Trois observations de cette lettre valent d'être retenues, parce que les pages
suivantes s'en servent sans y revenir.

D'abord, une remarque de régularité : si $A$ est un anneau local régulier,
$(x_{i})_{1 \leqslant i \leqslant r}$ une partie d'un système régulier de
paramètres et $(n_{i})$ des entiers $\geqslant 1$ \emph{non nécessairement
premiers à la caractéristique résiduelle}, alors
\[
  B \;=\; A[T_{1},\dots,T_{r}]\,/\,(T_{1}^{\,n_{1}} - x_{1},\ \dots,\
  T_{r}^{\,n_{r}} - x_{r})
\]
est encore régulier. C'est ce qui permet de ne pas regarder les
caractéristiques résiduelles des points non maximaux des $D_{i}$ : il en
résulte \emph{a posteriori} que si, en tout point maximal $s$ de chaque
$D_{i}$, le groupe d'inertie est d'ordre premier à la caractéristique de
$k(s)$, alors il l'est automatiquement en toute spécialisation de $s$ ---
énoncé vide sauf si $s$ est de caractéristique nulle.

Ensuite, une remarque de méthode : les démonstrations des paragraphes 7 et 8
« marchent pour n'importe quel type de données de ramification », pourvu qu'on
dispose d'un énoncé du type de son lemme 7.2. C'est la première indication que
la généralité coûte moins qu'elle ne rapporte, et c'est l'argument qu'il
opposera à Murre quand celui-ci proposera de se restreindre au cas kummérien.

Enfin, la lecture correcte des groupes de 10.2 et 10.3 : ce sont
$\pi_{1}(U,\;)$ et son plus grand quotient premier à $p$, et les théorèmes
devraient être énoncés en ces termes.

\subsection*{16 mai 1967 : Murre ne voit toujours pas}

La dernière lettre --- placée en tête du carton, écrite en dernier --- est la
plus utile pour lire l'esquisse, parce qu'elle pose deux questions dont les
réponses fixent des définitions que l'esquisse laisse incomplètes.

Murre demande si son intuition est juste : \emph{au lieu de donner directement
le revêtement $r(X)$ de $S$, on donne le comportement du revêtement au-dessus
d'autres revêtements, les $s(M)$, qui sont d'un type particulier} ; le
comportement au-dessus de tous les $s(M)$ ensemble en dit plus que $r(X)$
seul. C'est exactement cela, et c'est la meilleure phrase du lot pour dire ce
qu'est une donnée de ramification.

Il demande ensuite --- et c'est une lacune réelle --- si, dans la définition
d'un revêtement galoisien de type $\mathcal{R}$, on ne doit pas exiger que
l'action du groupe de Galois $G$ \emph{commute} à celle du groupe
d'automorphismes que les modèles $s(M)$ portent déjà. Grothendieck le confirme,
mais dans la lettre d'accompagnement, à la page 36 : la définition de la page
30 est incomplète telle qu'elle est écrite, et cette lecture-ci l'a complétée
là où elle se trouve.\footnote{La condition manquante est énoncée page 36 :
« the action of $G$ commutes to étale localization on $S$ \emph{and} to the
action of the groups of automorphisms given on the various $Z$'s ». Elle ne
figure pas au paragraphe 4 des pages 30 à 31, qui définit la notion. Le lot 2
de cette édition la porte à sa place.}

Il relève enfin, dans le cas kummérien, que la condition
$\underline{n}' \geqslant \underline{n}$ devrait s'écrire
$\underline{n}' = \underline{m} \cdot \underline{n}$. Grothendieck répond page
37 : « when I wrote $n' \geqslant n$, I meant of course the order relation of
divisibility ». La lecture ci-dessous écrit donc la divisibilité.

\pagerange{12}{14}
\section*{La donnée de ramification (pages 12 à 14)}

Soit $S$ un schéma et $\mathrm{Et}(S)$ son site étale. Une \emph{donnée de
ramification} sur $S$ est un couple
\[
  \mathcal{R} \;=\; (\mathcal{C},\, s)
\]
formé d'une catégorie fibrée $p : \mathcal{C} \to \mathrm{Et}(S)$ et d'un
foncteur cartésien
\[
  s : \mathcal{C} \longrightarrow \mathrm{Fl}_{S} ,
\]
où $\mathrm{Fl}_{S}$ est la catégorie fibrée sur $\mathrm{Et}(S)$ dont la fibre
en $S'$ est la catégorie des schémas relatifs sur $S'$.\footnote{Le tapuscrit
souligne les catégories : $\underline{C}$, $\underline{s}$, $\underline{Et}_S$,
$\underline{Fl}_S$, $\underline{R}$. Le soulignement est ici abandonné, comme
partout dans cette édition ; il ne portait aucune information que la casse ne
porte.} On ne suppose pas que $\mathcal{C}$ soit un \emph{champ} : ni les
objets ni les morphismes ne se recollent nécessairement. On ne suppose pas non
plus $s$ fidèle, encore moins pleinement fidèle --- et c'est le point de toute
l'affaire.

Un objet $M$ de la fibre $\mathcal{C}_{S'}$ s'appelle un \emph{modèle de
ramification}, et le $S'$-schéma $s(M)$ sa \emph{réalisation géométrique}. Dans
tous les exemples qui suivent, les $s(M)$ sont finis sur $S'$.

Composant $s$ avec le foncteur source, on obtient un foncteur
$\mathcal{C} \to (\mathrm{Sch})$, et l'on peut tirer en arrière le long de lui
la catégorie fibrée des schémas finis étales : on obtient une catégorie fibrée
$\mathcal{D}$ sur $\mathcal{C}$, dont la fibre en $M$ est la catégorie des
revêtements étales finis de $s(M)$.

\subsection*{Les revêtements de type $\mathcal{R}$}

Un \emph{crible couvrant} de $\mathcal{C}$ est une sous-catégorie pleine
$\mathcal{C}_{0}$ qui est un crible --- si l'extrémité d'une flèche est dans
$\mathcal{C}_{0}$, son origine y est --- et telle que les $p(M)$, pour $M$
parcourant $\mathcal{C}_{0}$, recouvrent $S$. Pour un tel crible, on pose
\[
  \mathrm{Rev}^{\mathcal{R}}_{\mathcal{C}_{0}}(S) \;=\;
  \text{catégorie des sections cartésiennes de } \mathcal{D}
  \text{ au-dessus de } \mathcal{C}_{0} ,
\]
et l'on définit la \emph{catégorie des revêtements de $S$ de type
$\mathcal{R}$} comme la limite inductive de ces catégories suivant les
restrictions :
\[
  \mathrm{Rev}^{\mathcal{R}}(S) \;=\;
  \varinjlim_{\mathcal{C}_{0}} \mathrm{Rev}^{\mathcal{R}}_{\mathcal{C}_{0}}(S) .
\]
Concrètement, un revêtement de type $\mathcal{R}$ est la donnée, pour chaque
modèle $M$ d'un crible couvrant, d'un revêtement étale fini $X(M)$ de $s(M)$,
et, pour chaque flèche $M' \to M$, d'un isomorphisme
$X(M') \simeq X(M) \times_{s(M)} s(M')$, ces isomorphismes se composant.
C'est la traduction exacte de la phrase de Murre : on ne donne pas le
revêtement, on donne son comportement au-dessus de tous les
modèles.\footnote{Le tapuscrit écrit tantôt
$\mathrm{Rev}^{\underline{R}}$, tantôt $\mathrm{Rev}_{\underline{R}}$, pour le
même objet, et parfois dans la même phrase (pages 13 et 14). L'exposant est
retenu ici pour tout le carton.}

Deux hypothèses accompagnent la définition, et il importe de les voir comme des
hypothèses.\footnote{Le tapuscrit écrit « Signalons que dans tous les cas que
nous aurons à utiliser, les foncteurs de transition \dots\ sont pleinement
fidèles », puis « Nous supposerons que \dots ». Ce sont bien deux conditions
posées sur les exemples, non deux propositions : la première est fausse pour
une donnée quelconque, et l'identification des
$\mathrm{Rev}^{\mathcal{R}}_{\mathcal{C}_{0}}(S)$ à des sous-catégories pleines
en dépend.}

\begin{enumerate}
  \item[(H1)] Les foncteurs de restriction
    $\mathrm{Rev}^{\mathcal{R}}_{\mathcal{C}_{0}}(S) \to
    \mathrm{Rev}^{\mathcal{R}}_{\mathcal{C}'_{0}}(S)$ sont pleinement fidèles.
    Alors les $\mathrm{Rev}^{\mathcal{R}}_{\mathcal{C}_{0}}(S)$ s'identifient,
    à équivalence près, à des sous-catégories pleines de
    $\mathrm{Rev}^{\mathcal{R}}(S)$, qui les épuise.
  \item[(H2)] Pour tout crible $\mathcal{C}_{0}$ et tout objet $X$ de
    $\mathrm{Rev}^{\mathcal{R}}_{\mathcal{C}_{0}}(S)$, le foncteur composé
    $\mathcal{C}_{0} \xrightarrow{\ X\ } \mathcal{D} \to (\mathrm{Sch})_{/S}$
    admet une limite inductive dans $(\mathrm{Sch})_{/S}$, notée $r(X)$, et
    cette limite ne change pas quand on restreint $X$ à un crible plus petit.
\end{enumerate}

Sous (H2) les $r$ partiels se raccordent en un unique foncteur
\[
  r : \mathrm{Rev}^{\mathcal{R}}(S) \longrightarrow (\mathrm{Sch})_{/S} ,
\]
la \emph{réalisation géométrique}. Les deux hypothèses sont vérifiées dans les
trois exemples ci-dessous, « de généralité croissante ».

Dans les cas qui intéressent, $r$ sera \emph{fidèle} sans être pleinement
fidèle. C'est là toute la mise : un revêtement de type $\mathcal{R}$ n'est pas
déterminé par sa réalisation géométrique, et il faut donc distinguer
soigneusement les deux. Quand $r$ est fidèle, on peut décrire
$\mathrm{Rev}^{\mathcal{R}}(S)$ comme la catégorie des « revêtements de $S$
munis d'une structure de type $\mathcal{R}$ », une telle structure sur un
revêtement $X$ étant la donnée d'une classe d'isomorphie d'objets
$X' \in \mathrm{Rev}^{\mathcal{R}}(S)$ et d'un isomorphisme
$r(X') \simeq X$.\footnote{La remarque est page 15, avec l'aveu qu'elle
« aurait dû figurer avant l'exemple ». Elle est remontée ici, où elle éclaire
la définition qu'elle commente.}

\pagerange{14}{16}
\section*{Exemple 1 : un revêtement à opérateurs (pages 14 à 16)}

C'est l'exemple à garder en tête, et Grothendieck le dira lui-même page 34 :
« the basic example in order to understand the idea is example 1 ».

Donnons-nous un $S$-schéma $Z$ muni d'une action à droite d'un groupe fini $G$
par $S$-automorphismes. Soit $\mathcal{C}_{0}$ la catégorie à un objet dont le
monoïde des endomorphismes est $G$, et prenons
$\mathcal{C} = \mathrm{Et}(S) \times \mathcal{C}_{0}$. Se donner un foncteur
cartésien de $\mathcal{C}$ dans une catégorie fibrée sur $\mathrm{Et}(S)$
revient à se donner un objet de la fibre au-dessus de $S$ muni d'une action de
$G$ ; le couple $(Z, G)$ fournit donc bien un $s$.

On trouve aussitôt une équivalence entre $\mathrm{Rev}^{\mathcal{R}}(S)$ et la
catégorie des revêtements étales finis $Z' \to Z$ munis d'une action à droite
de $G$ au-dessus de celle de $G$ sur $Z$ ; et la réalisation géométrique
existe si et seulement si le quotient $Z'/G$
existe, auquel cas $r(Z') = Z'/G$. Il suffit pour cela que toute orbite de $G$
dans $Z$ soit contenue dans un ouvert affine --- en particulier, il suffit que
$Z$ soit affine sur $S$, ce qui sera le cas puisque $Z$ sera un revêtement de
$S$.

\subsection*{Le sens intuitif}

En première approximation, les objets de $\mathrm{Rev}^{\mathcal{R}}(S)$ sont
les revêtements de $S$ qui \emph{deviennent étales} au-dessus du revêtement
ramifié donné $Z$. C'est ce qui donne son sens au mot « donnée de
ramification » : $Z$ est le patron de ramification que l'on s'autorise, et l'on
regarde tout ce qui n'est pas pire que lui.

Grothendieck le dit une seconde fois, page 34, dans le langage des corps de
fonctions, et cette formulation-là est la plus parlante. Supposons $S$ intègre
normal de corps de fonctions $K$, et $Z$ le normalisé de $S$ dans une extension
galoisienne finie $L$ de $K$. Un revêtement $S'$ de $S$, connexe, correspond à
une extension $K'$ de $K$. Dire que $S'$ n'a pas de ramification pire que
$Z/S$, c'est dire que l'image inverse normalisée de $S'$ au-dessus de $Z$ est
étale ; et lorsque $S$ est strictement local, cela revient tout simplement à
dire que \emph{$K'$ est isomorphe à une sous-extension de $L$}. La notion est
de nature locale pour la topologie étale, de sorte qu'on peut toujours se
ramener au cas strictement local, où elle est aussi simple que cela.

Alors, par transport de structure, $G$ opère sur $S'$, et l'on retrouve $S'$ à
partir du $G$-revêtement $Z'$ de $Z$ comme $S' = Z'/G$.

\subsection*{Quand $r$ est fidèle, quand il est pleinement fidèle}

Soit $U$ un ouvert de $S$ tel que $V = Z|U$ soit un revêtement principal de
groupe $G$, et supposons que toute partie ouverte et fermée de $Z$ stable par
$G$ et contenant $V$ soit égale à $Z$ --- c'est le cas si $V$ est dense dans
$Z$.

\begin{enumerate}
  \item[(a)] Alors $r$ est \emph{fidèle}. En effet le foncteur de restriction
    $Z' \mapsto Z'|V$, des $G$-revêtements étales de $Z$ dans ceux de $V$, est
    fidèle ; et le foncteur $Z'_{V} \mapsto Z'_{V}/G$, des $G$-revêtements
    étales de $V$ dans les revêtements étales de $U$, est une
    \emph{équivalence} par descente le long du revêtement principal
    $V \to U$ --- \emph{a fortiori} fidèle.
  \item[(b)] Si de plus $V$ est dense dans $Z$ et $Z$ est \emph{normal}, alors
    $Z' \mapsto Z'|V$ est pleinement fidèle, donc $r$ l'est aussi. Dans ce cas
    $\mathrm{Rev}^{\mathcal{R}}(S)$ est équivalente à une sous-catégorie
    \emph{pleine} de la catégorie des revêtements normaux de $S$ : ceux qui
    sont étales au-dessus de $U$ et dont toute composante irréductible domine
    une composante irréductible de $S$ rencontrant $U$.
\end{enumerate}

L'énoncé (b) est une équivalence sur une sous-catégorie pleine, non une
égalité : la sous-catégorie n'est pas décrite par une propriété fermée des
revêtements, et c'est précisément le théorème d'Abhyankar qui, dans le cas des
croisements normaux, la décrira ainsi.\footnote{Page 21 : « Le théorème
d'Abhyankar permet de caractériser ces revêtements par la condition d'être
étales sur $S - \operatorname{supp} D$, normaux en les points au-dessus de
$\operatorname{supp} D$ \dots ». C'est-à-dire : dans ce cas particulier, et
dans celui-là seulement, la structure devient une propriété. Voir le lot 2.}

Signalons enfin ce que la page 35 ajoute et qui n'est pas dans l'esquisse :
toutes les hypothèses de normalité sont \emph{superflues} pour la théorie
générale des revêtements étales de $(Z,G)$, qui a toutes les propriétés de
celle des revêtements étales d'un schéma, et est décrite par un groupe profini
--- extension de $G$ par $\pi_{1}(Z)$. Les hypothèses de normalité ne servent
qu'à assurer la pleine fidélité de $Z' \mapsto Z'/G$.

\pagerange{16}{18}
\section*{Exemple 2 : un système projectif de modèles (pages 16 à 18)}

Variante immédiate. Au lieu d'un seul couple $(Z, G)$, on se donne un système
projectif $(Z_{i}, G_{i})_{i \in I}$ de tels couples, $I$ ordonné filtrant. Le
système projectif des $G_{i}$ définit une catégorie fibrée $\mathcal{C}_{0}$
sur la catégorie associée à $I$, et l'on prend
$\mathcal{C} = \mathrm{Et}(S) \times \mathcal{C}_{0}$. Alors
\[
  \mathrm{Rev}^{\mathcal{R}}(S) \;\simeq\;
  \varinjlim_{i} \mathrm{Rev}^{\mathcal{R}_{i}}(S) ,
\]
limite inductive des catégories définies séparément par chaque $(Z_{i},G_{i})$.

Le cas intéressant est celui où les foncteurs de transition --- qui sont des
images inverses de revêtements étales à opérateurs --- sont pleinement fidèles,
de sorte que les $\mathrm{Rev}^{\mathcal{R}_{i}}(S)$ s'identifient à des
sous-catégories pleines qui épuisent la limite : c'est l'hypothèse (H1) sur cet
exemple. Il en est ainsi lorsque les morphismes de transition
$(Z_{i}, G_{i}) \to (Z_{j}, G_{j})$ sont tels que $G_{i} \to G_{j}$ soit
surjectif et que $Z_{i} \to Z_{j}$ induise un isomorphisme
$Z_{j} \simeq Z_{i}/N_{ij}$, $N_{ij}$ étant le noyau de $G_{i} \to G_{j}$. En
effet, tout revêtement étale $Z'_{j}$ de $Z_{j}$ vérifie alors
$Z'_{j} \simeq Z'_{i}/N_{ij}$, où $Z'_{i}$ est l'image inverse de $Z'_{j}$
sur $Z_{i}$.

La motivation est page 35 et elle est concrète : dans les applications on ne
dispose pas d'un seul modèle de ramification, mais d'une infinité --- tous les
$Z^{n}_{a}$ pour $a$ fixé et $n$ variable --- et la condition qu'on veut
imposer à un revêtement est d'avoir une ramification qui ne soit pire que celle
de \emph{l'un} des modèles.

\pagerange{18}{19}
\section*{Exemple 3, dans ses deux moutures (pages 18 et 19)}

L'exemple 3 est retapé, et le carton conserve les deux états. Ils ne sont pas
fusionnés ici : leur différence est la trace du travail, et elle va toujours
dans le même sens --- vers l'élimination d'une donnée arbitraire.

La situation qu'ils veulent couvrir est celle de la page 35 : les modèles de
ramification ne sont pas définis globalement --- dans le cas kummérien, les
diviseurs peuvent n'être pas principaux --- mais seulement localement, avec des
conditions de compatibilité disant que la restriction d'un modèle donné sur $U$
à $U \cap V$ n'a pas de ramification pire que celle d'un modèle donné sur $V$.

\subsection*{Première mouture : par une factorisation de $s$}

On suppose que $s$ se factorise
\[
  \mathcal{C} \xrightarrow{\ s'\ } \mathrm{Flop}_{S} \longrightarrow
  \mathrm{Fl}_{S} ,
\]
où $\mathrm{Flop}_{S}$ est la catégorie fibrée des schémas relatifs munis d'un
groupe fini d'opérateurs à droite, le second foncteur étant l'oubli des
opérations. Les morphismes de $\mathrm{Flop}_{S}$ sont les couples $(f, u)$ où
$f : Z \to Z'$ est un $S$-morphisme et $u : G_{S} \to G'_{S}$ un morphisme
\emph{localement constant, pas nécessairement constant}, de schémas en
groupes sur $S$, compatibles au sens évident.\footnote{Le tapuscrit écrit
« $f : X \to X'$ » là où ses propres objets sont $Z$ et $Z'$ ; c'est une faute
de frappe, non un changement d'objet, et elle est corrigée ici. La
transcription la conserve, comme elle doit.} La liberté laissée à $u$ est ce
qui permettra, au cas kummérien, de faire varier le groupe des racines de
l'unité d'un ouvert étale à l'autre.

On suppose alors deux choses. Pour toute flèche $M' \to M$ dans un
$\mathcal{C}_{S'}$, si $s(M') = (Z', G')$ et $s(M) = (Z,G)$, la flèche induit
un isomorphisme $Z \simeq Z'/H$, où $H$ est le noyau de $G'_{S} \to G_{S}$.
Et pour tout objet $M$, posant $s(M) = (Z,G)$, on se donne un homomorphisme
$G \to \operatorname{Aut}(M)$ tel que l'image $\bar{g}$ de $g$ vérifie
$s(\bar{g}) = (g_{Z}, \operatorname{int}(g))$, tout endomorphisme de $M$ étant
localement de cette forme ; le tout compatible aux images inverses et aux
flèches de $\mathcal{C}_{S'}$.

C'est cet homomorphisme $G \to \operatorname{Aut}(M)$ qui est la donnée
artificielle, et la seconde mouture s'en débarrasse.

\subsection*{Seconde mouture : par les groupes d'automorphismes}

On suppose maintenant, sur $\mathcal{R} = (\mathcal{C}, s)$ seule :

\begin{enumerate}
  \item[(1)] pour tout $S' \in \mathrm{Et}(S)$ et tout
    $M \in \mathcal{C}_{S'}$, le groupe $G = \operatorname{Aut}(M)$ est fini ;
  \item[(2)] pour toute flèche $M' \to M$ de $\mathcal{C}_{S'}$, il existe un
    unique morphisme de groupes $G' = \operatorname{Aut}(M') \to
    G = \operatorname{Aut}(M)$ rendant $M' \to M$ équivariant ; et si $H$ en est
    le noyau, $s(M') \to s(M)$ induit un isomorphisme
    $s(M')/H \simeq s(M)$ ;\footnote{Le sens de ce morphisme est fixé par deux
    flèches à l'encre portées au verso de la feuille (page 17) :
    $G' \to G$ au-dessus de $M' \to M$. Sans elles la phrase du recto est
    ambiguë. C'est le seul endroit du lot où un verso décide d'une lecture.}
  \item[(3)] deux morphismes de $M'$ dans $M$ se déduisent l'un de l'autre par
    composition avec un élément de $G$ ;
  \item[(4)] deux objets $M$, $M'$ de $\mathcal{C}_{S'}$ sont localement
    majorés : il existe une famille surjective de morphismes étales
    $S'_{i} \to S'$ telle que $M_{S'_{i}}$ et $M'_{S'_{i}}$ soient tous deux
    majorés par un même objet $M''_{i}$ de $\mathcal{C}_{S'_{i}}$.
\end{enumerate}

Sous ces conditions, la théorie de la descente montre que la réalisation
géométrique $r : \mathrm{Rev}^{\mathcal{R}}(S) \to (\mathrm{Sch})_{/S}$ est
définie, et à valeurs dans les revêtements de $S$ dès que les $s(M)$ en sont
--- il suffirait même, pour appliquer la descente, que les $s(M)$ soient
affines sur $S'$. Et les critères de fidélité et de pleine fidélité de
l'exemple 1 se transposent : il suffit que chaque
$\bigl(s(M), \operatorname{Aut}(M)\bigr)$ satisfasse, sur son $S'$, les
conditions que $(Z,G)$ y satisfaisait.

La condition (4) est celle qui fait le travail de recollement : c'est elle qui
dit que les modèles locaux, donnés sur des ouverts étales différents, sont
comparables. Le paragraphe des « Remarques » de la page 32 remplacera ces
conditions par cinq conditions a) à e) portant sur $\mathcal{R}$ seule, où la
donnée de $s$ elle-même a fondu ; c'est le troisième état de l'exemple 3, et il
est lu au lot 2.

\pagerange{19}{20}
\section*{Le cas kummérien (pages 19 et 20)}

Soit $(D_{i})_{i \in I}$ une famille de diviseurs de Cartier sur $S$. Pour
$S'$ étale sur $S$, on prend pour $\mathcal{C}_{S'}$ la catégorie dont les
objets sont les quadruples
\[
  M = (\underline{a},\, \underline{n},\, G,\, u) ,
\]
où $\underline{a} = (a_{i})_{i \in I}$ est une famille d'équations des
$D_{i\,S'}$, $\underline{n} = (n_{i})_{i \in I}$ une famille d'entiers
$\geqslant 1$ premiers aux caractéristiques résiduelles de $S'$, $G$ un groupe
fini, et $u : G_{S'} \xrightarrow{\ \sim\ } \mu_{\underline{n},\,S'}$ un
isomorphisme de schémas en groupes.

Les morphismes sont donnés par
\[
  \operatorname{Hom}_{S'}(M', M) \;=\;
  \begin{cases}
    \bigl\{\, \underline{b} = (b_{i})_{i \in I} \ :\
      \underline{b}^{\,\underline{n}} =
      \underline{a}'\,\underline{a}^{-1} \,\bigr\}
      & \text{si } \underline{n} \ \text{divise}\ \underline{n}' , \\[2pt]
    \varnothing & \text{sinon} ,
  \end{cases}
\]
la divisibilité s'entendant composante par composante.\footnote{Le tapuscrit
écrit « sauf si $\underline{n}' \geqslant \underline{n}$ ». Murre relève page 2
qu'il faudrait $\underline{n}' = \underline{m}\cdot\underline{n}$, et
Grothendieck répond page 37 : « when I wrote $n' \geqslant n$, I meant of
course the order relation of divisibility ». Le carton corrige donc lui-même,
à quarante pages de distance ; c'est la divisibilité qui est écrite ici.}

Le foncteur cartésien associe à $M$ le \emph{revêtement kummérien élémentaire}
\[
  Z^{\underline{n}}_{\underline{a}} \;=\;
  \operatorname{Spec} \mathcal{O}_{S'}[T_{i}\,;\ i \in I]\,/\,
  (T_{i}^{\,n_{i}} - a_{i}) ,
\]
considéré comme revêtement à groupe d'opérateurs $G$ grâce à
$u : G \simeq \mu_{\underline{n}}$. On vérifie qu'on est alors dans les
conditions de l'exemple 3.

Le foncteur $r$ est \emph{fidèle} ; il est \emph{pleinement fidèle} lorsque
$S$ est localement noethérien, normal aux points des
$\operatorname{supp} D_{i}$, les $D_{i}$ réduits, et toute composante
irréductible d'un $D_{i}$ distincte de toute composante d'un $D_{j}$ pour
$j \neq i$.

Le groupe $G$ et l'isomorphisme $u$ sont ici encore des données arbitraires :
Murre demande page 2 si les objets ne sont pas plutôt les couples
$(\underline{a}, \underline{n})$ avec $G(M) = \mu_{\underline{n}}$, et
Grothendieck le confirme page 37. La présentation intrinsèque de la page 33
permet en effet de « se dispenser d'introduire aussi des isomorphismes
$u : G_{S'} \to \mu_{\underline{n},\,S'}$ », et l'objet se réduit au couple
$(\underline{a}, \underline{n})$. C'est l'état auquel le carton reviendra vingt
fois ; sous cette forme, la catégorie fibrée $\mathcal{C}$ est un groupoïde,
et son champ associé est une \emph{gerbe} sous $\mu_{\underline{n}}$ --- ce que
la page 200 dira explicitement, et ce qu'on appelle aujourd'hui le champ des
racines $\sqrt[\underline{n}]{\underline{D}/S}$.\footnote{Aucun de ces mots
n'est disponible en 1967 : les gerbes sont de Giraud (1971), et les champs des
racines de Cadman et d'Abramovich--Graber--Vistoli (années 2000). L'objet, lui,
est ici, avec sa classe d'obstruction dans $H^{2}$ ; voir la lecture du lot 10,
qui l'écrit là où Grothendieck l'écrit.}

Le lot 2 continue le tapuscrit à la page 21, avec le second cas kummérien
--- celui d'un diviseur $D$ à croisements normaux dont les composantes ne sont
définies que localement pour la topologie étale --- et c'est là que le mot
« modérément ramifié au sens large » est prononcé.

\end{document}
